\section{Discussion and Future Work}

\label{sec:discussion}

\subsection{Implications for Agent Design}
\label{sec:implications}

The Harness-Hacker Protocol reframes a fundamental assumption of current agent
design. The harness is not immutable infrastructure; it is one of the most
valuable targets for agent improvement. By treating self-modification as a
first-class capability with calibrated safety controls, HHP unlocks a class of
compound improvements that purely behavioral adaptation cannot achieve.

The practical implication is architectural: agents should be designed with
self-modification hooks from the start. The \texttt{PluginAPI} surface, the
skill library schema, the worktree lifecycle management---these are not
afterthoughts but core infrastructure. An agent built with HHP in mind will
outperform an equivalent agent without it, because it will converge to the
optimal tool set for its deployment context rather than being limited to the
tools its designers anticipated.

The \emph{specific-tools-beat-generic-flexibility} principle ($93.1\%$ vs.\
$83.8\%$) quantifies the performance cost of the current paradigm. Every
percentage point of that gap represents real task failures that HHP can
systematically eliminate.

\subsection{Automated PR Review by a Second Agent}
\label{sec:auto-review}

The current Tier 4 gate requires human PR review. This is appropriate for
production systems but introduces latency: the agent must wait for a human
reviewer before the improvement can be deployed. An important extension is
automated review by a second, independent agent instance.

The reviewing agent would receive the diff, the CI results, the originating
friction event, and the agent's justification. It would evaluate: (a) is the
modification narrowly scoped to address the identified friction? (b) are there
plausible security or correctness risks that CI did not catch? (c) is the
proposed implementation the simplest adequate solution?

Two-agent review is weaker than human review (both agents may share systematic
biases) but stronger than no review, and it would dramatically reduce the
deployment latency for non-security-sensitive modifications. We propose a
tiered review policy: automated review for Tier 4 modifications not touching
security-sensitive paths, human review for those that do.

\subsection{Cross-Fleet Propagation via DSO}
\label{sec:dso-propagation}

Successful harness modifications are, in effect, structured experiences that
should be shared across the agent fleet. The Decentralized Semantic Optimization
protocol~\cite{hanzo2026dso} provides exactly the infrastructure needed:
Byzantine-robust aggregation of experiential priors across distributed agent
populations.

A Tier 2 extension that proves highly effective on one agent's deployments
should be propagated to other agents. The DSO gossip protocol can carry
extension metadata (name, trigger pattern, performance delta, trust score) to
nodes that have not yet deployed the extension. Nodes can pull and evaluate
extensions that match their own telemetry profiles, applying the same trust
escalation logic locally.

This creates a collective evolution dynamic: successful harness improvements
spread through the fleet via DSO, and the fleet as a whole converges toward
the optimal tool set faster than any individual agent could achieve alone.

\subsection{Formal Verification of Generated Tools}
\label{sec:verification}

The current validation gate $\Gamma_1$ tests generated tools against a finite
set of triggering failure cases. This is sound for the specific cases tested
but does not provide guarantees for novel inputs. A natural extension is to
apply lightweight formal verification---property-based testing~\cite{claessen2000quickcheck}
or automated theorem proving for simple properties---to generated tools before
hot-loading.

For tools with well-typed interfaces, type-level guarantees (TypeScript strict
mode, runtime schema validation) provide a first layer of formal assurance.
For tools with more complex correctness properties, symbolic execution or
model checking could verify invariants beyond what unit tests cover.

\subsection{Hierarchical Scope Control}
\label{sec:scope-control}

The current ClassifyScope algorithm (Algorithm~\ref{alg:classify}) makes
tier decisions based on simple heuristics. A more principled approach would
use a learned classifier trained on historical friction events and their
outcomes---using the GRPO-extracted telemetry to build a meta-policy over
self-modification decisions. This meta-policy would learn which friction types
are best addressed at which tier, improving the efficiency of the modification
budget.

% ============================================================================
% 10. CONCLUSION
% ============================================================================
